\documentclass[CJK,14pt]{beamer}
\usepackage{CJKutf8}
\usepackage{beamerthemesplit}
\usetheme{Malmoe}
\useoutertheme[footline=authortitle]{miniframes}
\usepackage{amsmath}
\usepackage{amssymb}
\usepackage{graphicx}
\usepackage{eufrak}
\usepackage{color}
\usepackage{slashed}
\usepackage{simplewick}
\usepackage{tikz}
\usepackage{tcolorbox}
\usepackage{ulem}
\graphicspath{{../figures/}}
%%figures
\def\lfig#1#2{\includegraphics[width=#1 in]{#2}}
\def\tfig#1#2{\includegraphics[height=#1 in]{#2}}
\def\addfig#1#2{\begin{center}\includegraphics[width=#1 in]{#2}\end{center}}
\def\wulian{\includegraphics[width=0.18in]{emoji_wulian.jpg}}
\def\bigwulian{\includegraphics[width=0.35in]{emoji_wulian.jpg}}
\def\bye{\includegraphics[width=0.18in]{emoji_bye.jpg}}
\def\bigbye{\includegraphics[width=0.35in]{emoji_bye.jpg}}
\def\huaixiao{\includegraphics[width=0.18in]{emoji_huaixiao.jpg}}
\def\bighuaixiao{\includegraphics[width=0.35in]{emoji_huaixiao.jpg}}
\def\jianxiao{\includegraphics[width=0.18in]{emoji_jianxiao.jpg}}
\def\bigjianxiao{\includegraphics[width=0.35in]{emoji_jianxiao.jpg}}
\def\haoqi{\includegraphics[width=0.18in]{emoji_haoqi.jpg}}
%% colors
\def\blacktext#1{{\color{black}#1}}
\def\bluetext#1{{\color{blue}#1}}
\def\redtext#1{{\color{red}#1}}
\def\darkbluetext#1{{\color[rgb]{0,0.2,0.6}#1}}
\def\skybluetext#1{{\color[rgb]{0.2,0.7,1.}#1}}
\def\cyantext#1{{\color[rgb]{0.,0.5,0.5}#1}}
\def\greentext#1{{\color[rgb]{0,0.7,0.1}#1}}
\def\darkgray{\color[rgb]{0.2,0.2,0.2}}
\def\lightgray{\color[rgb]{0.6,0.6,0.6}}
\def\gray{\color[rgb]{0.4,0.4,0.4}}
\def\blue{\color{blue}}
\def\red{\color{red}}
\def\orange{\color[rgb]{1.,0.8,0.}}
\def\green{\color{green}}
\def\darkgreen{\color[rgb]{0,0.4,0.1}}
\def\darkblue{\color[rgb]{0,0.2,0.6}}
\def\skyblue{\color[rgb]{0.2,0.7,1.}}
%%control
\def\uline#1{\underline{\ \hspace{#1 in}\ }}
\def\bral{\left(\begin{array}{l}}
\def\brar{\end{array}\right)}
\def\brall{\left(\begin{array}{ll}}
\def\brarr{\end{array}\right)}
\def\bralll{\left(\begin{array}{lll}}
\def\brarrr{\end{array}\right)}
\def\branchl{\left\{\begin{array}{l}}
\def\branchr{\end{array}\right.}
\def\branchll{\left\{\begin{array}{ll}}
\def\branchrr{\end{array}\right.}
\def\branchlll{\left\{\begin{array}{lll}}
\def\branchrrr{\end{array}\right.}
\def\sfgamma#1{\,\Gamma\left( #1 \right)\,}
\def\be{\begin{equation}}
\def\ee{\nonumber\end{equation}}
\def\bea{\begin{eqnarray}}
\def\eea{\nonumber\end{eqnarray}}
\def\bch{\begin{CJK}{UTF8}{gbsn}}
\def\ech{\end{CJK}}
\def\bchl{\bch\begin{large}}
\def\echl{\end{large}\ech}
\def\bchL{\bch\begin{Large}}
\def\echL{\end{Large}\ech}
\def\bitem{\begin{itemize}}
\def\eitem{\end{itemize}}
\def\bcenter{\begin{center}}
\def\ecenter{\end{center}}
\def\bex{\begin{minipage}{0.2\textwidth}\includegraphics[width=0.6in]{jugelizi.png}\end{minipage}\begin{minipage}{0.76\textwidth}}
\def\eex{\end{minipage}}
\def\chtitle#1{\frametitle{\bch#1\ech}}
\def\bmat#1{\left(\begin{array}{#1}}
\def\emat{\end{array}\right)}
\def\bcase#1{\left\{\begin{array}{#1}}
\def\ecase{\end{array}\right.}
\def\bmini#1{\begin{minipage}{#1\textwidth}}
\def\emini{\end{minipage}}
\def\tbox#1{\begin{tcolorbox}#1\end{tcolorbox}}
\def\pfrac#1#2#3{\left(\frac{\partial #1}{\partial #2}\right)_{#3}}
\def\res#1#2{\,\mathrm{res}\,#1\left(#2\right)\,}
\def\newt#1#2{\left(\begin{array}{c}#1\\ #2\end{array}\right)}
\def\reof#1{\,\mathrm{Re}{\left(#1\right)\,}}
\def\imof#1{\,\mathrm{Im}{\left(#1\right)\,}}
\def\Arg#1{\,\mathrm{Arg}\,#1\,}
%%symbols
%%symbols
\def\bropt{\,(\ \ \ )}
\def\cvmol{C_V^{\rm mol}}
\def\cpmol{C_p^{\rm mol}}
\def\intfull{\int_{-\infty}^\infty}
\def\inthalf{\int_0^\infty}
\def\bropt{\,(\ \ \ )}
\def\optlist#1#2#3{(A)\;#1\ \ \ (B)\;#2\ \ \ (C)\;#3}
\def\ooptlist#1#2#3#4{(A)\;#1\ \ \ (B)\;#2\ \ \ (C)\;#3\ \ (D)\;#4}
\def\sone{$\star$}
\def\stwo{$\star\star$}
\def\sthree{$\star\star\star$}
\def\sfour{$\star\star\star\star$}
\def\sfive{$\star\star\star\star\star$}
\def\rint{{\int_\leftrightarrow}}
\def\roint{{\oint_\leftrightarrow}}
\def\stdHf{{\textit{\r H}_f}}
\def\deltaH{{\Delta \textit{\r H}}}
\def\ii{{\dot{\imath}}}
\def\skipline{{\vskip0.1in}}
\def\skiplines{{\vskip0.2in}}
\def\lagr{{\mathcal{L}}}
\def\hamil{{\mathcal{H}}}
\def\vecv{{\mathbf{v}}}
\def\vecx{{\mathbf{x}}}
\def\vecy{{\mathbf{y}}}
\def\veck{{\mathbf{k}}}
\def\vecp{{\mathbf{p}}}
\def\vecn{{\mathbf{n}}}
\def\vecA{{\mathbf{A}}}
\def\vecP{{\mathbf{P}}}
\def\vecsigma{{\mathbf{\sigma}}}
\def\hatJn{{\hat{J_\vecn}}}
\def\hatJx{{\hat{J_x}}}
\def\hatJy{{\hat{J_y}}}
\def\hatJz{{\hat{J_z}}}
\def\hatj#1{\hat{J_{#1}}}
\def\hatphi{{\hat{\phi}}}
\def\hatq{{\hat{q}}}
\def\hatpi{{\hat{\pi}}}
\def\vel{\upsilon}
\def\Dint{{\mathcal{D}}}
\def\adag{{\hat{a}^\dagger}}
\def\bdag{{\hat{b}^\dagger}}
\def\cdag{{\hat{c}^\dagger}}
\def\ddag{{\hat{d}^\dagger}}
\def\hata{{\hat{a}}}
\def\hatb{{\hat{b}}}
\def\hatc{{\hat{c}}}
\def\hatd{{\hat{d}}}
\def\hatD{{\,\hat{D}}}
\def\hatN{{\hat{N}}}
\def\hatH{{\hat{H}}}
\def\hatp{{\hat{p}}}
\def\Fup{{F^{\mu\nu}}}
\def\Fdown{{F_{\mu\nu}}}
\def\newl{\nonumber \\}
\def\vece{\mathrm{e}}
\def\calM{{\mathcal{M}}}
\def\calT{{\mathcal{T}}}
\def\calR{{\mathcal{R}}}
\def\barpsi{\bar{\psi}}
\def\baru{\bar{u}}
\def\barv{\bar{\upsilon}}
\def\qeq{\stackrel{?}{=}}
\def\ftf{\stackrel{\mathcal{FT}}{\Longrightarrow}}
\def\ftb{\stackrel{\mathcal{FT}}{\Longleftarrow}}
\def\ftfb{\stackrel{\mathcal{FT}}{\Longleftrightarrow}}
\def\ltf{\stackrel{\mathcal{LT}}{\Longrightarrow}}
\def\ltb{\stackrel{\mathcal{LT}}{\Longleftarrow}}
\def\ltfb{\stackrel{\mathcal{LT}}{\Longleftrightarrow}}
\def\torder#1{\mathcal{T}\left(#1\right)}
\def\rorder#1{\mathcal{R}\left(#1\right)}
\def\contr#1#2{\contraction{}{#1}{}{#2}#1#2}
\def\trof#1{\mathrm{Tr}\left(#1\right)}
\def\trace{\mathrm{Tr}}
\def\comm#1{\ \ \ \left(\mathrm{used}\ #1\right)}
\def\tcomm#1{\ \ \ (\text{#1})}
\def\slp{\slashed{p}}
\def\slk{\slashed{k}}
\def\calp{{\mathfrak{p}}}
\def\veccalp{\mathbf{\mathfrak{p}}}
\def\Tthree{T_{\tiny \textcircled{3}}}
\def\pthree{p_{\tiny \textcircled{3}}}
\def\dbar{{\,\mathchar'26\mkern-12mu d}}
\def\erf{\mathrm{erf}}
\def\const{\mathrm{const.}}
\def\pheat{\pfrac p{\ln T}V}
\def\vheat{\pfrac V{\ln T}p}
%%units
\def\fdeg{{^\circ \mathrm{F}}}
\def\cdeg{^\circ \mathrm{C}}
\def\atm{\,\mathrm{atm}}
\def\angstrom{\,\text{\AA}}
\def\SIL{\,\mathrm{L}}
\def\SIkm{\,\mathrm{km}}
\def\SIyr{\,\mathrm{yr}}
\def\SIGyr{\,\mathrm{Gyr}}
\def\SIV{\,\mathrm{V}}
\def\SImV{\,\mathrm{mV}}
\def\SIeV{\,\mathrm{eV}}
\def\SIkeV{\,\mathrm{keV}}
\def\SIMeV{\,\mathrm{MeV}}
\def\SIGeV{\,\mathrm{GeV}}
\def\SIcal{\,\mathrm{cal}}
\def\SIkcal{\,\mathrm{kcal}}
\def\SImol{\,\mathrm{mol}}
\def\SIN{\,\mathrm{N}}
\def\SIHz{\,\mathrm{Hz}}
\def\SIm{\,\mathrm{m}}
\def\SIcm{\,\mathrm{cm}}
\def\SIfm{\,\mathrm{fm}}
\def\SImm{\,\mathrm{mm}}
\def\SInm{\,\mathrm{nm}}
\def\SImum{\,\mathrm{\mu m}}
\def\SIJ{\,\mathrm{J}}
\def\SIW{\,\mathrm{W}}
\def\SIkJ{\,\mathrm{kJ}}
\def\SIs{\,\mathrm{s}}
\def\SIkg{\,\mathrm{kg}}
\def\SIg{\,\mathrm{g}}
\def\SIK{\,\mathrm{K}}
\def\SImmHg{\,\mathrm{mmHg}}
\def\SIPa{\,\mathrm{Pa}}
%page
\def\secpage#1#2{\begin{frame}\bch\bcenter{\bf \Huge #1} \skipline \tbox{#2}\ecenter\ech\end{frame}}

\begin{document}
\bch

\title{第四课：最短路径问题}
\subtitle{http://zhiqihuang.top/algorithms}
  \author{黄志琦}
  \date{}

  \maketitle

  \begin{frame}
    \frametitle{内容摘要}
  \tableofcontents
  \end{frame}

  \section{图的最短路径问题}

  \begin{frame}
    \frametitle{寒假快到了，苗十三准备去旅游了！}


    因为这学期被{\bf 扣工资扣得很惨}，苗十三手头不是很宽裕，必须精打细算。他想弄清去各个城市的最低开销。

    {\scriptsize（嗯，不用考虑回来的开销。苗十三准备找警察叔叔说自己被拐卖，免费被送回来。）}


  \end{frame}


  \begin{frame}
    \frametitle{苗十三查到的各个航班的机票价格}
    \addfig{4}{paths.jpg}
  \end{frame}
  

  \begin{frame}
    \frametitle{用矩阵来表示}
    用0,1,2,\ldots,7分别表示珠海，上海，北京，广州，重庆，南京，杭州，拉萨。

    \skiplines
    
    用一个二维数组 x[8][8] 来表示航班价格：

    \skiplines
    
    x[i][j] = i到j的直飞价格（如无航班记作$\infty$)

  \end{frame}

  \section{Floyd-Warshall算法}

  \begin{frame}
    \frametitle{Floyd-Warshall算法思想}

    第一步：把矩阵更新为

    x[i][j] = 允许到城市0中转，i到j的最少机票花销

    {\blue \small
  if(x[i][0]+x[0][j]$<$x[i][j]) x[i][j]=x[i][0]+x[0][j];}

\skipline

第二步：把矩阵更新为

x[i][j] = 允许到城市0和1中转，i到j的最少机票花销

{\blue \small
  if(x[i][1]+x[1][j]$<$x[i][j]) x[i][j]=x[i][1]+x[1][j];}
    
{\scriptsize 注意：如果最佳转机顺序是i$\rightarrow$1$\rightarrow$0$\rightarrow$j.因为x[1][j]已经在第一步被更新为 i$\rightarrow$0$\rightarrow$j的价格，在第二步仍然可以通过比较x[i][1]+x[1][j]和x[i][j]找到正确的最低价格。}

这样一直进行下去，第八步之后(搞定)：

x[i][j] = 允许任意转机, i到j的的最少机票花销
  \end{frame}
  
  \begin{frame}
    \frametitle{Floyd-Warshall算法的优缺点}
    \bitem
  \item{优点：算法及其简明；可以求出图中任意两点最短距离。}
  \item{缺点：对$N$个顶点的图，算法复杂度是$O(N^3)$，不适合太多个顶点的问题。}
    \eitem
  \end{frame}

  \begin{frame}
    \frametitle{程序实现}

    参考：
    
    zhiqihuang.top/algorithms/examples.php
    
    
  \end{frame}

  \section{Dijkstra算法}
  
  \begin{frame}
    Floyd-Warshall算法一下子把任何两个城市间的最低机票开销算出来了。但苗十三只关心从珠海到各个城市怎么飞最便宜。有没有不那么“浪费”的算法呢？

    \skiplines
    
    有的，比如Dijkstra算法就是个不错的方法。
  \end{frame}
  

  \begin{frame}
    \frametitle{Dijkstra算法的基本思想}
    用一个额外的数组m来记录珠海到各个城市的最少机票开销：

    如果要求直飞：

    m=\{$600,900,200,1100,\infty,\infty,\infty$\};

      
    
  \end{frame}


  \begin{frame}
    \frametitle{Dijkstra算法的基本思想}
    注意到珠海到广州的直飞价格200最低。显然，即使允许任意转机，这个价格也是最低的（因为一往其他地方飞，就超过这个价格了)。我们用蓝色标记出来，表示这个价格已经是允许任意转机的最低价了：

    m=\{$600,900,{\blue 200},1100,\infty,\infty,\infty$\};      

    接下来，我们立刻考虑能否通过从广州转一次机省钱：从广州可以直接到达的只有北京和拉萨，经过比较发现可以把到北京和拉萨的最低开销都更新：

    m = \{ $600, 800, {\blue 200}, 1100, \infty, \infty, 1700$ \};

      
    
  \end{frame}


  \begin{frame}
    \frametitle{Dijkstra算法的基本思想}
    现在还没有标记蓝色的数字中最小的是珠海到上海的价格600。这个其实也已经是允许任意转机的最低价了！为什么呢？因为要比这个价格低，只可能从价格更低的广州转机。而从广州转机的价格我们已经考虑并更新过了。

    OK，既然是允许任意转机的最低价，我们继续标记蓝色：
    

    m = \{ ${\blue 600}, 800, {\blue 200}, 1100, \infty, \infty, 1700$ \};

    可以从上海转机去重庆和南京，我们立刻比较更新一下去重庆和南京的最低开销：

    m = \{ ${\blue 600}, 800, {\blue 200}, 1000, 800, \infty, 1700$ \};    
  \end{frame}


\begin{frame}
  \frametitle{Dijkstra算法的基本思想}
  现在的
  
  m = \{ ${\blue 600}, 800, {\blue 200}, 1000, 800, \infty, 1700$ \};

  记录的是允许从广州和上海转机的最低开销。还没蓝染（em？）的城市中最低开销有北京800和南京800。这两个也都已经是允许任意转机的最低开销了！都蓝染一下：
  
  m = \{ ${\blue 600}, {\blue 800}, {\blue 200}, 1000, {\blue 800}, \infty, 1700$ \};  

  可以从北京和南京转机去上海（已经蓝了别指望更新了），杭州，拉萨；比较后发现可以把去杭州最低开销更新一下：

    m = \{ ${\blue 600}, {\blue 800}, {\blue 200}, 1000, {\blue 800}, 1100, 1700$ \};  
  \end{frame}
    

\begin{frame}
  \frametitle{Dijkstra算法的基本思想}
  继续找黑色的最小那个数，重庆：1000。标上蓝色：

  m = \{ ${\blue 600}, {\blue 800}, {\blue 200}, {\blue 1000}, {\blue 800}, 1100, 1700$ \};

  可以从重庆转机去的南京已经是蓝色，不需要考虑更新。
\end{frame}


\begin{frame}
  \frametitle{Dijkstra算法的基本思想}
  继续找黑色的最小那个数，杭州：1100。标上蓝色：

  m = \{ ${\blue 600}, {\blue 800}, {\blue 200}, {\blue 1000}, {\blue 800}, {\blue 1100}, 1700$ \};

  可以从杭州转机去的上海已蓝，跳过。
\end{frame}


\begin{frame}
  \frametitle{Dijkstra算法的基本思想}
  最后黑色的拉萨：1700.标上蓝色：

  m = \{ ${\blue 600}, {\blue 800}, {\blue 200}, {\blue 1000}, {\blue 800}, {\blue 1100}, {\blue 1700}$ \};

  最后的答案就是这些全部染蓝的数字！
\end{frame}


\begin{frame}
  \frametitle{Dijkstra算法的复杂度}
  假如每个顶点出去的边数不是很多的话，因为每次可以用二分法搜索“黑色最小数”，复杂度就是$O(N\log N)$。

  假如每个顶点出去的边数非常多，在最坏的情况下，就会是$O(N^2\log N)$的复杂度；当然，实际应用中这样的情形极少出现。

\end{frame}

\begin{frame}
  \frametitle{作业}
  用Dijkstra算法写程序解决苗十三的旅游问题。
\end{frame}



\begin{frame}
  \frametitle{补充说明}
  我们之前讨论的问题都是边的权重为正的情形。如果存在负权重边（就是航空公司倒贴苗十三请求苗十三飞），就要用一种叫Bellman-Ford的算法。因为应用场景不多，就不介绍了。
\end{frame}

  \ech
\end{document}
